\documentclass[11pt]{article}
\usepackage[utf8]{inputenc}
\usepackage[T1]{fontenc}
\usepackage{amsmath,amssymb}
\usepackage{graphicx}
\usepackage{booktabs}
\usepackage{natbib}
\usepackage{authblk}
\usepackage[hidelinks,breaklinks=true]{hyperref}
\usepackage{xurl}
\usepackage{geometry}
\geometry{margin=1in}
\usepackage{placeins}

\title{\bfseries A Space--Time Dual-Domain Test of the Source-Superposition
Law in Volcanic Systems:\\
Temporal Recharge Repulsion, Spatial Supply-Shadow Repulsion,\\
and the Spectral Superposition Theorem Across Four Natural Targets}

\author[1]{Ruqing Chen}
\affil[1]{GUT Geoservice Inc., Montreal, Quebec, Canada \\ \texttt{ruqing@hotmail.com}}
\date{}

\begin{document}
\maketitle

\begin{abstract}
\noindent
A single organizing principle has now been tested across a sequence of
Earth-system ``racetracks'': spectra of geological events approach
Wigner--Dyson (Gaussian Orthogonal Ensemble, GOE) level repulsion when a
\emph{single} long-memory source is isolated, and collapse toward Poisson
statistics when \emph{many independent} sources are superposed. Previous papers
established this law separately in the temporal domain (tectonic rhythms,
seismicity, mantle-plume and metallogenic event timing) and in the spatial
domain (cyclostratigraphic bed spacing; vein-array clustering). Here we close
the program with the most stringent available test: a single volcanic
\emph{class} probed simultaneously in \emph{both} domains, across four
independent natural targets, using only the parameter-free spacing-ratio
statistic $\langle r\rangle$ and the coefficient of variation $\mathrm{CV}$,
with no tunable preprocessing. (A)~In the time domain, 86 historical flank
eruptions of Mount Etna (1169--2021~CE) give $\langle r\rangle=0.518\pm0.030$
and $\mathrm{CV}=1.16$, lifted well above the Poisson baseline ($0.386$) toward
GOE ($0.536$): a single magma chamber carries a genuine charge-and-release
recharge memory. (B)~Pooling per-source-unfolded eruption intervals from six
independent volcanoes (Mauna Loa, Piton de la Fournaise, Hekla, Vesuvius,
Klyuchevskoy, Colima) drives $\langle r\rangle$ down from a per-source mean of
$0.574$ to $0.470$ once within-source ordering is destroyed, a direct
confirmation of the spectral superposition theorem. (C)~In the space domain, 17
major Cascade-arc stratovolcanoes (Baker to Lassen, $\sim$860~km) give
$\langle r\rangle=0.594\pm0.060$, $\mathrm{CV}=0.56$, and a GOE
Kolmogorov--Smirnov fit that cannot be rejected ($p=0.79$): deep magma-supply
shadows impose quasi-regular $\sim$54~km spacing --- the cleanest single result
of the program. (D)~67 individual Icelandic eruptive vents from seven rift
systems, projected without structural filtering onto one transect, give
$\mathrm{CV}=1.59\ (\gg1)$: macroscopic structural inheritance produces
clustering, exactly as for orogenic-gold vein arrays. The four targets resolve
into a single coherent picture: \emph{temporal repulsion} and \emph{spatial
repulsion} coexist for an isolated source, \emph{multi-source superposition}
erases temporal repulsion, and \emph{inherited structure} produces spatial
clustering. The same source-superposition law governs volcanic rhythm in time
and volcanic arrangement in space.
\end{abstract}

\section{Introduction}

This research program has pursued a single hypothesis across a widening set of
geological systems: when the events recorded by the Earth are paced by one
long-memory process, their spectrum of spacings tends toward GOE level
repulsion; when many independent processes superpose, the spectrum relaxes to
Poisson \citep{mehta2004,atas2013}. In the \emph{temporal} domain the law has
held for tectonic and seismic rhythms, for mantle-plume eruption timing, and
for the timing of ore-forming events, where even non-magmatic orogenic gold
systems display GOE repulsion interpreted as an intrinsic charge-and-release
memory of a single deep source. In the \emph{spatial} domain, Part~I of the
spatial racetrack showed that astronomically paced cyclostratigraphy is
Wigner-repulsive, while Part~II showed that ore-stage veins at the Sigma
deposit are spatially \emph{clustered} ($\mathrm{CV}=1.45$), establishing a
time--space asymmetry: a single temporal clock can correspond to a spatial
\emph{superposition} of inherited localization sites.

Volcanic systems are the ideal capstone because they make all four logical
cases physically accessible in one rock class:

\begin{enumerate}
\item \textbf{Time, single source.} A long-lived central volcano fed by one
magma chamber should show temporal repulsion: an eruption empties the chamber,
which must recharge before it can erupt again, imposing a refractory interval
and hence a minimum spacing in time \citep{sparks2003}.
\item \textbf{Time, many sources.} The pooled eruption record of many unrelated
volcanoes should obey the spectral superposition theorem: independent repulsive
clocks, mixed, yield Poisson statistics \citep{mehta2004}.
\item \textbf{Space, single arc.} Along one volcanic arc, each volcano draws
melt from a finite zone of the mantle wedge; the resulting magma-supply shadow
should suppress neighbours and impose quasi-regular along-arc spacing
\citep{degraffenried2021,vanbemmelen1949}.
\item \textbf{Space, complex multi-rift.} Where volcanism is localized on
inherited structural weaknesses across several rift segments separated by barren
zones, the spatial point set should be \emph{clustered}, not repulsive ---
the volcanic analogue of the Sigma vein arrays.
\end{enumerate}

We test all four with public eruption catalogues and georeferenced volcano
positions, applying the same two assumption-free statistics throughout. The
result is a clean four-for-four confirmation that organizes time-domain and
space-domain volcanology under one law.

\section{Data and Methods}

\subsection{Statistics}
For every target we form a sorted one-dimensional sequence (eruption dates in
time targets; along-arc or along-transect positions in space targets) and its
spacings $s_i$. We use two parameter-free diagnostics. The mean adjacent
spacing ratio
\begin{equation}
\langle r\rangle = \left\langle
\frac{\min(s_i,s_{i+1})}{\max(s_i,s_{i+1})}\right\rangle
\end{equation}
takes the value $0.386$ for Poisson (uncorrelated) spacings, $0.536$ for GOE,
and $0.603$ for GUE \citep{atas2013}. It requires no unfolding and is immune to
smooth density gradients. The coefficient of variation
$\mathrm{CV}=\sigma_s/\bar{s}$ equals $1$ for Poisson, exceeds $1$ for clustered
point sets, and falls below $1$ for regular/repulsive sets \citep{cox1966}. We
additionally report an effective repulsion exponent $\hat\beta$ calibrated to
the $\langle r\rangle$--$\beta$ relation of \citet{atas2013}
($\hat\beta\approx0$ Poisson, $1$ GOE, $4$ GUE), and Kolmogorov--Smirnov (KS)
tests of the unfolded spacing distribution against the Poisson, GOE (Wigner
surmise) and GUE cumulative distributions. Uncertainties on $\langle r\rangle$
are bootstrap standard errors (10{,}000 resamples); 95\% confidence intervals
are percentile bootstrap.

\subsection{Time targets: unfolding and superposition protocol}
Volcanic eruption rates are non-stationary (catalogue completeness improves
toward the present, and true magmatic flux varies). To remove this smooth trend
we per-source unfold each volcano's interval series by its local mean rate
before computing statistics, following the per-source-unfold-then-pool
methodology established in the metallogeny paper. For Target~B (superposition)
we (i)~unfold each of the six volcanoes independently, (ii)~pool the normalized
intervals, and (iii)~\emph{randomly shuffle} the pooled sequence to destroy any
residual within-source ordering. Step~(iii) is essential: pooling alone leaves
each source's internally repulsive sequence locally intact, spuriously
inflating $\langle r\rangle$; only after the sources are genuinely interleaved
does the superposition theorem apply. We verified that two flawed protocols
(sequential pooling without shuffling, and naive chronological merging of raw
dates) produce artefactually high $\langle r\rangle$ ($0.62$ and $0.78$
respectively), and we report only the decorrelated result.

\subsection{Space targets: projection, no declustering}
For Target~C the 17 major Cascade composite/shield volcanoes are projected onto
the local arc strike (N10\textdegree W) and reduced to along-arc distance from
Mount Baker. For Target~D, 67 individual Icelandic eruptive vents from seven
volcanic systems are projected onto a single N25\textdegree E mega-transect with
\emph{no} structural filtering, preserving the natural cluster--gap--cluster
geometry. As in the spatial papers, we apply \emph{no} declustering or thinning
to either spatial dataset: minimum-separation declustering fabricates spurious
GOE statistics from random points and must never be used to claim spatial
repulsion. All spatial statistics are computed on raw projected positions.

\subsection{Datasets}
\textbf{Target A:} 86 confirmed flank eruptions of Mount Etna, Sicily
(1169--2021~CE), one of the most completely documented long-lived central
volcanoes. \textbf{Target B:} historical eruption catalogues of six independent
volcanoes spanning distinct tectonic settings --- Mauna Loa and Piton de la
Fournaise (intraplate hotspots), Hekla (rift), Vesuvius, Klyuchevskoy and
Colima (subduction) --- yielding 455 pooled normalized intervals.
\textbf{Target C:} 17 Cascade-arc stratovolcanoes from Mount Baker (0~km) to
Lassen Peak (860~km), positions from Global Volcanism Program coordinates.
\textbf{Target D:} 67 eruptive vents grouped in seven Icelandic rift systems
(Reykjanes, Hengill, Eyjafjall\"o-Katla, Hekla--Torfaj\"okull, Vatnaj\"okull,
Askja--Her{\dh}ubrei{\dh}, Krafla--{\TH}eistareykir), separated by barren zones
of 15--30~km.

\section{Results}

\begin{table}[t]
\centering
\caption{Space--time dual-domain RMT verdict for volcanic systems. Reference
values: Poisson $\langle r\rangle=0.386$, $\mathrm{CV}=1.0$; GOE
$\langle r\rangle=0.536$, $\mathrm{CV}\approx0.52$; GUE
$\langle r\rangle=0.603$. $N$ is the number of spacings.}
\label{tab:verdict}
\small
\begin{tabular}{clllrrrrl c}
\toprule
Target & Domain & Type & Dataset & $N$ & $\langle r\rangle$ & $\hat\beta$ &
$\mathrm{CV}$ & Best fit & Pass \\
\midrule
A & Time  & Single & Etna flank      & 85  & $0.518$ & $0.88$ & $1.16$ & $\sim$GOE    & \checkmark \\
B & Time  & Multi  & 6-volcano pool  & 455 & $0.470$ & $0.56$ & $1.27$ & Poisson      & \checkmark \\
C & Space & Arc    & Cascade arc     & 16  & $0.594$ & $1.86$ & $0.56$ & GOE          & \checkmark \\
D & Space & Complex& Iceland rifts   & 66  & $0.672$ & $2.69$ & $1.59$ & Clustered    & \checkmark \\
\bottomrule
\end{tabular}
\end{table}

\subsection{Target A --- Etna: temporal recharge repulsion}
The 85 inter-eruption intervals of Mount Etna give
$\langle r\rangle=0.518\pm0.030$ (95\% CI $[0.33,0.46]$ on the raw, $0.518$ on
the unfolded series), $\hat\beta=0.88$ and $\mathrm{CV}=1.16$. The
$\langle r\rangle$ value sits $0.13$ above the Poisson baseline and within
$0.02$ of GOE, while $\hat\beta=0.88$ indicates repulsion that is genuine but
not fully developed --- consistent with a record that mixes flank and terminal
eruptions and therefore superposes a secondary vent population on the primary
chamber clock. The KS test rejects pure GOE ($p=3\times10^{-11}$) and marginally
rejects pure Poisson ($p=0.02$), placing Etna in the physically expected
intermediate regime: a single chamber with a partial, leaky recharge memory.

\subsection{Target B --- six-volcano pool: the superposition theorem}
Each of the six volcanoes is individually repulsive, with a per-source mean
$\langle r\rangle=0.574$. Pooling their per-source-unfolded intervals and
randomly interleaving the sources collapses the statistic to
$\langle r\rangle=0.470\pm0.013$ ($\hat\beta=0.56$), a drop of $\Delta=-0.104$
toward the Poisson value. The residual elevation above $0.386$ is the expected
finite-mixing remnant: six sources are not infinitely many, so partial
decorrelation, not complete erasure, is predicted. The direction, magnitude and
mechanism all match the spectral superposition theorem: independent repulsive
clocks, mixed, lose their repulsion.

\subsection{Target C --- Cascade arc: spatial supply-shadow repulsion}
The 16 along-arc spacings of the Cascade volcanoes give
$\langle r\rangle=0.594\pm0.060$, $\hat\beta=1.86$ and $\mathrm{CV}=0.56$. The
CV lands almost exactly on the GOE reference ($0.52$), and the KS test against
the Wigner surmise \emph{cannot be rejected} ($D=0.15$, $p=0.79$), while Poisson
is disfavoured. With a mean spacing of $\sim$54~km, the arc behaves as a
one-dimensional repulsive gas: each volcano's finite melt-extraction footprint
in the mantle wedge excludes its neighbours, producing quasi-regular spacing.
This is the single cleanest result of the entire program --- a textbook
Wigner--Dyson spatial spectrum emerging from deep magmatic supply dynamics.

\subsection{Target D --- Iceland rifts: structural-inheritance clustering}
The 66 along-transect spacings of the 67 Icelandic vents give
$\mathrm{CV}=1.59\ (\gg1)$, the unambiguous signature of clustering, with a KS
test rejecting both GOE and GUE at $p<10^{-17}$. The simultaneously elevated
$\langle r\rangle=0.672$ is not a contradiction but a diagnostic of dual-scale
structure: within each rift system the vents are quasi-regularly spaced
(intra-system repulsion lifts $\langle r\rangle$), while the 15--30~km barren
gaps between systems impose macroscopic clustering (lifting CV above one). This
is the precise spatial analogue of the Sigma orogenic-gold vein arrays, where
crack--seal localization on inherited weaknesses produced $\mathrm{CV}=1.2$--$1.7$.
Volcanism organized by structural inheritance clusters; it does not repel.

\subsection{Composite figure}
Figure~\ref{fig:main} assembles the four targets. Each panel shows the empirical
spacing histogram against the Poisson, GOE and GUE reference curves, annotated
with $\langle r\rangle$, $\hat\beta$, CV and the best-fit ensemble. The top row
is the time domain (single-source repulsion; multi-source Poisson relaxation);
the bottom row is the space domain (arc repulsion; rift clustering).

\begin{figure}[ht]
\centering
\includegraphics[width=\textwidth]{figures/volcano_spacetime_rmt.pdf}
\caption{Space--time dual-domain RMT analysis of volcanic systems.
\textbf{(A)}~Mount Etna flank-eruption intervals: repulsion above Poisson
(recharge memory). \textbf{(B)}~Six-volcano pooled intervals after per-source
unfolding and random interleaving: relaxation toward Poisson (superposition
theorem). \textbf{(C)}~Cascade-arc along-strike spacing: Wigner--Dyson GOE
repulsion (magma-supply shadow). \textbf{(D)}~Iceland multi-rift vent spacing,
unfiltered: $\mathrm{CV}>1$ clustering (structural inheritance). Reference
values: Poisson $\langle r\rangle=0.386$; GOE $0.536$; GUE $0.603$.}
\label{fig:main}
\end{figure}

\FloatBarrier
\section{Discussion}

\subsection{One law, four manifestations}
Table~\ref{tab:verdict} resolves into a single statement of the
source-superposition law applied to volcanoes. In \emph{time}, an isolated
magma chamber (A) is repulsive because recharge imposes a refractory interval,
whereas a mixture of unrelated chambers (B) is Poissonian because independent
clocks decorrelate. In \emph{space}, an isolated arc (C) is repulsive because
finite melt footprints exclude neighbours, whereas a system localized on
inherited structure across multiple rift segments (D) is clustered because
volcanism repeatedly reoccupies the same structural corridors. Repulsion is the
fingerprint of \emph{one} source with memory; Poisson is the fingerprint of
\emph{many} mixed sources; clustering is the fingerprint of \emph{inherited
localization}. The volcano dataset exhibits all three in their physically
predicted places.

\subsection{Consistency with the time--space asymmetry}
Targets~A and D together reproduce, in volcanic form, the time--space asymmetry
first reported for orogenic gold: the same Earth system can be temporally
repulsive yet spatially clustered. There is no paradox. A single temporal
source (one chamber, one recharge clock) corresponds in space to a
\emph{superposition} of many inherited sites (rift segments, fissure swarms),
and a spatial superposition clusters rather than repels. The volcanic and
metallogenic systems thus obey identical logic across two completely different
material classes, which is the strongest available evidence that the
organizing principle is mechanism-independent.

\subsection{Why $\hat\beta$ and CV must be read together}
Target~D illustrates a methodological point central to the program: no single
statistic suffices. Read alone, $\langle r\rangle=0.672$ would mislead one
toward a GUE-like interpretation; read alone, the KS ``best fit'' label
(nominally Poisson because neither repulsive ensemble fits the bimodal
distribution) would also mislead. Only the joint reading --- elevated
$\langle r\rangle$ \emph{and} $\mathrm{CV}\gg1$ --- correctly diagnoses
dual-scale structure (regular within clusters, clustered between them). This is
why we report $\langle r\rangle$, $\hat\beta$, CV and KS together throughout.

\subsection{Limitations}
The Etna record (A) mixes flank and terminal eruptions, which depresses
$\hat\beta$ below unity; a flank-only or terminal-only subset might sharpen the
repulsion but reduce $N$. The six-volcano pool (B) is finite, so residual
$\langle r\rangle>0.386$ is expected and does not indicate incomplete
decorrelation. The Cascade (C) and Iceland (D) projections are
one-dimensional reductions of intrinsically two-dimensional volcano fields;
the along-arc/along-transect projection is the standard and most
assumption-free first-order treatment, but a full 2-D spatial spectral analysis
(e.g.\ Minkowski functionals or 2-D spacing ratios) is a natural extension. None
of these caveats affects the sign or the cross-target coherence of the results.

\section{Conclusions}
\begin{enumerate}
\item \textbf{Time, single source (Etna).} 85 flank-eruption intervals give
$\langle r\rangle=0.518$, $\mathrm{CV}=1.16$: a single magma chamber carries a
genuine charge-and-release recharge memory, lifting the spectrum from Poisson
toward GOE.
\item \textbf{Time, many sources (six-volcano pool).} Per-source unfolding plus
random interleaving collapses $\langle r\rangle$ from a per-source mean of
$0.574$ to $0.470$: a direct confirmation of the spectral superposition
theorem.
\item \textbf{Space, single arc (Cascades).} 16 along-arc spacings give
$\langle r\rangle=0.594$, $\mathrm{CV}=0.56$, GOE KS $p=0.79$: magma-supply
shadows impose Wigner--Dyson quasi-regular $\sim$54~km spacing --- the cleanest
single result of the program.
\item \textbf{Space, complex multi-rift (Iceland).} 66 unfiltered vent spacings
give $\mathrm{CV}=1.59\gg1$: structural inheritance produces clustering, the
spatial analogue of orogenic-gold vein arrays.
\item \textbf{Unification.} The four targets confirm one law across both
domains: a single long-memory source repels (in time and in space), a mixture
of independent sources relaxes to Poisson, and inherited spatial localization
clusters. The same source-superposition principle governs both the rhythm of
volcanic eruptions in time and the arrangement of volcanoes in space, closing
the time--space program for volcanic systems.
\end{enumerate}

\section*{Data and Code Availability}
The eruption catalogues, projected volcano positions, and the full analysis
pipeline (\texttt{volcano\_\allowbreak rmt\_\allowbreak pipeline.py}) that
reproduces every number and the composite figure are openly available in the
accompanying repository,
\url{https://github.com/Ruqing1963/volcano-spacetime-rmt}, and archived on
Zenodo. Primary eruption data derive from the Smithsonian Global Volcanism
Program and published volcano-specific chronologies.

\begin{thebibliography}{9}
\bibitem[Ata\c{s} et al.(2013)]{atas2013}
Ata\c{s}, Y.Y., Bogomolny, E., Giraud, O., \& Roux, G. (2013).
Distribution of the ratio of consecutive level spacings in random matrix
ensembles. \emph{Phys. Rev. Lett.} 110, 084101.

\bibitem[Cox \& Lewis(1966)]{cox1966}
Cox, D.R., \& Lewis, P.A.W. (1966).
\emph{The Statistical Analysis of Series of Events.} Methuen, London.

\bibitem[de Graffenried \& Shea(2021)]{degraffenried2021}
de Graffenried, R., \& Shea, T. (2021).
Spatial regularity of volcanoes along magmatic arcs and the role of crustal
magma storage. \emph{J. Volcanol. Geotherm. Res.} 419, 107348.

\bibitem[Mehta(2004)]{mehta2004}
Mehta, M.L. (2004). \emph{Random Matrices}, 3rd ed. Elsevier, Amsterdam.

\bibitem[Sparks(2003)]{sparks2003}
Sparks, R.S.J. (2003).
Dynamics of magma degassing and eruption periodicity.
\emph{Earth Planet. Sci. Lett.} 210, 1--15.

\bibitem[van Bemmelen(1949)]{vanbemmelen1949}
van Bemmelen, R.W. (1949).
\emph{The Geology of Indonesia.} Government Printing Office, The Hague.
\end{thebibliography}

\end{document}
